\begin{exsol}{25}{Cabral 3.20 Productivity (applied)}{64}{Advanced}{A2}
\asked Obtain firm-level data on outputs and inputs (quantities and prices) and estimate each
firm's total factor productivity with the method of \S3.1 (Equation 3.2).

\method TFP as the multiplicative coefficient of a Cobb-Douglas technology common to the
industry, recovered as a residual, with the exponents measured by cost shares (Cabral \S3.1,
``Productivity'' and ``Estimating total factor productivity'', pp.~39--40, eqs.~3.1--3.2;
Exercise 3.15 for the share result). The benchmark for dispersion is the evidence in \S4.2
(p.~75): within US four-digit industries the 90th-percentile plant has close to twice the TFP of
the 10th, and close to five times in China and India. Course link: ``organization matters''
predicts large, persistent productivity gaps between firms with the same inputs (\S3.5,
pp.~56--59; \MR{} Ch.~1).

\solution
\subparagraph{(a) The method.} Assume, as Cabral does (p.~39), that (i) all firms face the same
output price, (ii) each firm $i$ has $q_i = \omega_i K_i^{\alpha} L_i^{\beta}$, and (iii)
$\alpha$ and $\beta$ are common while $\omega_i$ differs. Taking logs,
\[
\ln \omega_i = \ln q_i - \alpha \ln K_i - \beta \ln L_i . \tag{3.2}
\]
A price-taking, profit-maximizing firm sets $\alpha\,p q_i = r K_i$ and $\beta\,p q_i = w L_i$
(the first-order conditions of $\max\ p\,\omega_i K^\alpha L^\beta - rK - wL$), so the exponents
equal the revenue shares of capital and labour costs:
\[
\hat\alpha = \overline{\left(\frac{rK}{pq}\right)}, \qquad
\hat\beta = \overline{\left(\frac{wL}{pq}\right)} ,
\]
averaged across firms. Plug $\hat\alpha$, $\hat\beta$ into (3.2) to get each $\ln\hat\omega_i$,
normalize by the industry mean, and summarize the dispersion by the standard deviation of
$\ln\hat\omega$ and the 90/10 ratio $\exp(p_{90} - p_{10})$.

\subparagraph{(b) Measurement issues.}
\begin{enumerate}[label=(\roman*)]
  \item \emph{Labour productivity is the wrong measure.} Revenue per worker rises with capital
  even when two firms share the same technology (p.~39). Under profit maximization it can say
  almost nothing: $pq/L = w/\beta$, so it reflects the wage, not efficiency.
  \item \emph{Quantity versus revenue.} Most datasets report revenue. Dividing by an industry
  price index measures quantity only if all firms charge the same price (assumption (i)). With
  firm-specific prices, ``revenue TFP'' mixes efficiency with market power and demand, the
  productivity-versus-profitability distinction of Foster, Haltiwanger and Syverson (2008), which
  Cabral cites in Chapter~4.
  \item \emph{Shares and market power.} With a markup, revenue shares understate the exponents;
  using shares of total cost (under constant returns) is the usual correction.
  \item \emph{Capital and its user cost} are hard to measure (book values, depreciation, the
  rate $r$).
  \item \emph{Regression instead of shares.} Estimating $\alpha$, $\beta$ by OLS of $\ln q$ on
  $\ln K$, $\ln L$ is biased, because firms choose inputs knowing their own $\omega_i$: more
  productive firms use more inputs, so OLS credits inputs with part of the productivity. Olley
  and Pakes (1996), also cited by Cabral, address this.
\end{enumerate}

\subparagraph{(c) A worked run on SYNTHETIC data.} No real dataset is shipped with these notes,
so the script \texttt{tfp\_example.py} (in this folder) demonstrates the method on
\textbf{simulated firms}. It draws 500 price-taking, profit-maximizing firms with true
$\alpha = 0.30$, $\beta = 0.50$ and log-TFP standard deviation $0.27$ (chosen so that the true
90/10 ratio is about 2, the US benchmark of p.~75), firm-specific input prices, small
input-choice errors and output measurement error (seed 2017). \textbf{The numbers below describe
these simulated firms and are not evidence about any real industry.} Output of
\texttt{python tfp\_example.py}:
{\small
\begin{verbatim}
SYNTHETIC DATA (seed 2017), 500 simulated firms; true alpha=0.30, beta=0.50
cost-share estimates: alpha=0.300  beta=0.502
log TFP: sd=0.274   90/10 TFP ratio=2.02
naive OLS of ln q on ln K, ln L: alpha=0.491  beta=0.495
corr(log labour productivity, log TFP) = 0.057
true log TFP: sd=0.276   90/10 ratio=2.01   corr(estimated, true)=0.983
\end{verbatim}
}
Reading: the cost-share method recovers the exponents and the TFP ranking (correlation $0.98$
with the true values) and the dispersion (90/10 ratio $2.02$ against $2.01$). The naive OLS
regression overstates $\alpha$ ($0.49$ against $0.30$) because input choices respond to
$\omega_i$, illustrating issue~(v). Labour productivity is nearly uncorrelated with TFP
($0.06$), illustrating issue~(i).

\subparagraph{(d) Using real data.} Save a CSV with one row per firm (same industry and year) and
columns \texttt{q, K, L, p, r, w}: output (quantity, or revenue deflated by an industry price
index), capital, labour, output price, user cost of capital and wage. Then run
\texttt{python tfp\_example.py firms.csv}; the same estimates are printed, labelled as CSV
data. Typical sources are manufacturing censuses and annual industrial surveys with
establishment microdata (for Brazil, IBGE's PIA through a secure data room) or commercial
balance-sheet panels; with revenue data, issue~(ii) applies and the result is revenue TFP.

\interp On real data the method typically finds what \S4.2 reports: large dispersion (about 2:1
between the 90th and 10th percentile within narrow US industries) that persists over time.
Firms with the same inputs produce very different outputs. That residual is where the course
puts organization: management practices, incentive systems and the fit among them (\S3.5,
p.~58; \MR{} Ch.~1). TFP measures the size of the gap; it does not by itself show that
organization causes it, which is why Cabral turns to field experiments for causality (p.~58).

\verify (1) With the true $\alpha$, $\beta$, the script's estimates are within $0.002$ of the
truth, and the script asserts $|\hat\alpha - 0.3| < 0.02$, $|\hat\beta - 0.5| < 0.02$ and a
correlation above $0.9$. (2) Limit case: with no TFP dispersion the residual would contain only
measurement and input-choice error (s.d.\ about $0.06$ in the simulation, which alone gives a 90/10 ratio of about
$e^{2 \times 1.2816 \times 0.058} \approx 1.16$), so a measured 90/10 ratio near 2
cannot come from noise of that size. (3) Analytic check of issue~(i): under the first-order
conditions $pq/L = w/\beta$ exactly, so the small correlation in the output is the expected
one, not a coding error. (4) Running the script on a CSV written from the synthetic firms
reproduces the same estimates, confirming the CSV path.

\answer{Estimate $\alpha$, $\beta$ as average revenue shares of capital and labour costs, then
$\ln\hat\omega_i = \ln q_i - \hat\alpha\ln K_i - \hat\beta\ln L_i$ (eq.~3.2), and report the
spread (s.d., 90/10 ratio). Watch for revenue versus quantity, markups, capital measurement and
the simultaneity bias of OLS. On \emph{synthetic} data the script recovers $\hat\alpha = 0.300$,
$\hat\beta = 0.502$ and a 90/10 ratio of $2.02$ (true $2.01$). Real data typically show ratios
near 2 (US) to 5 (China, India), the dispersion that ``organization matters'' predicts.}
\end{exsol}
